%% Elsevier CAS 风格初稿夹具（测试用最小骨架）
%%
%% 钉住四类真实踩过的坑：
%%   1. 跨行的 \newcommand —— 逐行搬运会截成半句，留下未闭合的 `{`，
%%      把整个文件的后续解析全带崩；
%%   2. `\author[a]{Bo Liu}\cormark[1]\ead{a@b.c}` —— 邮箱与通讯标记都写在
%%      `\author{}` **外面**，只读参数就两个字段全丢；
%%   3. `\affiliation[a]{organization={...}, city={...}, country={...}}` ——
%%      key=value 写法不能按逗号切（organization 的值自身含逗号）；
%%   4. 包里混着官方 kit 的其他主文件、说明书、缩略图（LTX050）——
%%      这些「正文一个字都没提」的文件不该进结果包。

\documentclass[a4paper,fleqn]{cas-sc}

\usepackage[numbers,sort&compress]{natbib}
\usepackage{graphicx}
\usepackage{amsmath,amssymb,amsfonts}
\usepackage{tikz}
\usepackage{float}

% 自定义命令（**跨行定义** —— 这一条是重点）
\newcommand*\circled[1]{\tikz[baseline=(char.base)]{
    \node[shape=circle, draw, inner sep=1pt] (char) {#1};}}
\renewcommand{\algorithmicrequire}{\textbf{Input:}}
\renewcommand{\algorithmicensure}{\textbf{Output:}}

\begin{document}

%标题
\title [mode = title]{Explainability-Guided Partial Privileged Learning: A Class-Aware Weighted Similarity Approach}
\shorttitle{}
\shortauthors{}

\author[a]{Bo Liu}\cormark[1]\ead{csbliu@gdut.edu.cn}\credit{Writing - original draft}
\author[b]{Yanshan Xiao}\ead{xiaoyanshan@gdut.edu.cn}\credit{Funding acquisition}

\affiliation[a]{organization={School of Automation, Guangdong University of Technology},
                city={Guangzhou},
                country={China}}
\affiliation[b]{organization={School of Computer Science, Guangdong University of Technology},
                city={Guangzhou},
                country={China}}

\cortext[cor1]{Corresponding author}
\fntext[1]{}

\begin{keywords}
    Learning Using Privileged Information (LUPI) \sep Partial Privileged Learning \sep Explainable AI (XAI) \sep SVM+
\end{keywords}

\maketitle

\begin{abstract}
Learning Using Privileged Information (LUPI) enhances classifier performance by
leveraging additional knowledge during training. We propose a framework termed
Explainability-Guided Partial Privileged Learning (EPPL) that imputes missing
privileged information with a class-aware weighted similarity strategy.
\end{abstract}

% Main text
\section{Introduction}\label{introduction}

Supervised learning has become a dominant paradigm in machine learning \cite{vapnik2009new}.
In practical scenarios we often have access to a priori information associated
with the training samples, which provides richer context for model
interpretation \cite{li2018r}. When a subset of samples lacks this privileged
component, standard LUPI methods cannot be directly applied.

\section{Methodology}\label{sec:methodology}

We denote the training set by $\{(x_i, y_i)\}$ and the privileged features by
$x_i^{*}$. The weighted cosine similarity between a query sample and a
candidate sample is computed with respect to the class-conditional importance
vector. Features with high discriminative power retain their weights, while
irrelevant or noisy features are strictly suppressed \cite{li2018r}.

\begin{figure}[t]
\centering
\includegraphics[width=0.8\linewidth]{fig1.png}
\caption{Class-aware weighted similarity.}\label{fig:similarity}
\end{figure}

% 留档备查的旧图（**注释掉的引用不算引用**，不该被带进结果包）
% \includegraphics{unused_fig.png}

Figure~\ref{fig:similarity} shows the weighted similarity computation.

\section{Conclusion}\label{conclusion}

We presented an explainability-guided framework for partial privileged
learning. Future work will extend the framework to multi-modal settings.

\section*{Declaration of competing interest}

The authors declare that they have no known competing financial interests or
personal relationships that could have appeared to influence the work reported
in this paper.

% 原文档类专有命令（定义在 cas-common.sty 里，Springer 的 sn-jnl.cls 没有）——
% 留在正文里编译必挂，转换时应注释保留并记 CVT030。
\printcredits

\bibliographystyle{cas-model2-names}
\bibliography{references}

\end{document}
